\chapter{Fluoroscopy}

\section*{What Is This Test?}
Fluoroscopy is an imaging technique that uses X-rays to obtain real-time moving images of the interior of an object. It acts like an 'X-ray movie' to guide procedures or evaluate dynamic bodily functions.

\section*{Alternative Names}
Continuous X-ray, Live X-ray

\section*{Principle}
Fluoroscopy is an imaging technique that uses X-rays to obtain real-time moving images of the interior of an object. It acts like an 'X-ray movie' to guide procedures or evaluate dynamic bodily functions.

\section*{History}
Thomas Edison invented the first commercial fluoroscope in 1896. The technology evolved from darkroom viewing to modern digital image intensifiers and flat-panel detectors.

\section*{Why This Test Is Ordered}
Ordered to guide catheter placement, visualize gastrointestinal tract function (barium swallow), guide joint injections, or evaluate cardiac pacemakers.

\section*{Is This a Primary or Secondary Test}
Primary guidance modality for interventional radiology and gastroenterology procedures.

\section*{How This Test Is Performed}
You lie or stand between an X-ray source and a detector. A continuous beam of X-rays is passed through the body part, displaying live motion on a monitor.

\section*{What Preparation Is Needed}
Preparation varies by examination. Follow department instructions; fasting is required only for selected studies or sedation. Remove objects that could interfere with imaging.

\section*{Contraindications and Precautions}
Pregnancy (unless absolutely necessary and shielded). Contrast allergy if contrast agents are used.

\section*{Risks}
Higher radiation exposure than standard X-rays due to continuous imaging. Skin burns are possible with extremely prolonged exposure.

\section*{What Happens After the Test}
If barium was used, drink plenty of water to prevent constipation. Resume normal activities.

\section*{Understanding Your Results}
Live movement shows normal transit of contrast, proper catheter placement, or normal joint mechanics. Obstructions or leaks are visible immediately.

\section*{Normal Results}
Normal physiological movement and transit of contrast; no leakage, obstruction, or mechanical failure.

\section*{Follow-up Tests}
Specific clinical follow-up depending on the procedure guided (e.g., surgical planning, endoscopy).

\section*{References}
\begin{enumerate}
  \item MedlinePlus. Fluoroscopy. U.S. National Library of Medicine; 2025.
  \item Current Medical Diagnosis and Treatment. McGraw-Hill; 2025.
\end{enumerate}

\clearpage
